\documentclass{aiaa-tc}

\usepackage{amsmath}
\usepackage{amssymb}
\usepackage{bm}
\usepackage{booktabs}
\usepackage{graphicx}
\usepackage{url}

\graphicspath{{figures/}}

\newcommand{\R}{\mathbb{R}}
\newcommand{\SO}{\mathrm{SO}}
\newcommand{\SE}{\mathrm{SE}}
\newcommand{\SEtwo}{\mathrm{SE}_2(3)}
\newcommand{\Exp}{\operatorname{Exp}}
\newcommand{\Log}{\operatorname{Log}}
\newcommand{\diag}{\operatorname{diag}}
\newcommand{\N}{\mathcal{N}}
\newcommand{\V}{\mathcal{V}}
\newcommand{\T}{\mathsf{T}}
\newcommand{\dd}{\mathrm{d}}

\title{Particle Filtering with Lie-Group Camera Projection for Online Multi-Camera UAV Tracking}

\author{Siddharth Anand\thanks{M.S. Student, School of Aeronautics and Astronautics, Purdue University, West Lafayette, IN.}}

\begin{document}

\maketitle

\begin{abstract}
This paper studies online tracking of a small unmanned aerial vehicle using
multiple calibrated external cameras.  The implemented estimator compares a
MEKF-style error-state extended Kalman filter, an iterated EKF, their gated
variants, and a bootstrap particle filter.  All filters use the same
Lie-group camera model: a world point is transformed by a calibrated
\(\SE(3)\) camera extrinsic and projected through a pinhole intrinsic matrix.
Monte Carlo simulations in a Gazebo/Purdue-world style camera geometry show
that the EKF variants are fastest and most accurate in the nearly unimodal
nominal cases, while the particle filter is more appropriate when the visual
posterior is multimodal.  A candidate-set experiment demonstrates this regime:
the particle filter achieves \(0.186\) m RMSE, compared with \(0.319\) m for
the best gated EKF baseline, because it marginalizes over multiple plausible
detections instead of committing to a single association.  A hardware-calibration
hypothesis test is also performed by replaying an extracted PURT motion in
simulation while injecting camera intrinsic and extrinsic errors.  The ablation
shows that using a mismatched \(1600\times 1200\) intrinsic matrix on
\(640\times 480\) images, or treating an active-marker body frame as an optical
camera frame, can reproduce multi-meter particle-filter errors.  Finally, a
future extension is formulated in which each particle carries a latent
\(\SE_2(3)\) inertial-navigation state so that body-frame acceleration readings
enter the prediction step through the \(\SO(3)\) exponential map and group
action, while the reported output remains the UAV trajectory.
\end{abstract}

\section*{Nomenclature}

\begin{tabular}{ll}
\(p_{w,k}\) & UAV position in the world frame at time index \(k\) \\
\(v_{w,k}\) & UAV velocity in the world frame at time index \(k\) \\
\(z_{j,k}\) & Pixel measurement from camera \(j\) at time index \(k\) \\
\(K_j\) & Intrinsic matrix of camera \(j\) \\
\(T_{j,cw}\) & World-to-camera extrinsic transform for camera \(j\), an element of \(\SE(3)\) \\
\(R_{j,cw},t_{j,cw}\) & Rotation and translation blocks of \(T_{j,cw}\) \\
\(\pi(\cdot)\) & Pinhole camera projection map \\
\(N_p\) & Number of particles in the particle filter \\
\(w_k^{(i)}\) & Weight of particle \(i\) at time index \(k\) \\
\(\omega_{m,k}\) & Body-frame gyroscope measurement \\
\(a_{m,k}\) & Body-frame accelerometer measurement or specific force \\
\(R_{wb,k}\) & Body-to-world attitude of the UAV \\
\(\V_k\) & Set of cameras with valid detections at time index \(k\)
\end{tabular}

\section{Introduction}

External vision sensors are attractive for tracking small UAVs in laboratory,
urban-canyon, and GPS-denied settings, including Purdue urban-canyon style
external-vision work \cite{yang2022}.  A multi-camera system can provide
multiple simultaneous bearing-like measurements, but its accuracy depends on
camera intrinsics, camera-to-world extrinsics, image timing, and detector
quality.  The estimator must also handle missing views, noisy detections, and
ambiguous visual associations.

This project investigates such a system in the context of online UAV trajectory
estimation.  The current implemented state is position and velocity,
\[
x_k =
\begin{bmatrix}
p_{w,k}\\
v_{w,k}
\end{bmatrix}
\in \R^6,
\]
while the camera model is written using \(\SE(3)\) transformations.  The
implemented filters are a MEKF-style error-state EKF, an iterated EKF (ItEKF),
gated versions of both EKF variants, and a bootstrap particle filter (PF).
Particle filters are emphasized because they can represent non-Gaussian and
multimodal posteriors \cite{thrun2005}, which are natural in camera-based
tracking when detections are ambiguous.

The paper has four objectives.  First, it presents the common Lie-group
projection model used by all filters.  Second, it compares the EKF-style
estimators and PF in controlled Monte Carlo simulations.  Third, it identifies
a regime in which PF is the more appropriate estimator: ambiguous candidate-set
measurements.  Fourth, it reproduces a hardware failure hypothesis in
simulation by injecting camera calibration errors observed in a PURT hardware
data collection.  A proposed extension then moves more Lie-group structure into
the PF prediction step by carrying latent attitude and using body-frame
acceleration readings.

\section{Problem Formulation}

\subsection{Online Multi-Camera Tracking Task}

At each time \(k\), the estimator receives detections from up to \(N_c\)
calibrated cameras.  The objective of the implemented system is to estimate
\(p_{w,k}\), with \(v_{w,k}\) carried to improve prediction between camera
updates.  The camera detections are two-dimensional pixel measurements; the
filters do not receive ground-truth 3D position during online operation.

The controlled simulation uses a \(12\) s trajectory with time step
\(\Delta t=0.05\) s, \(N_p=450\) particles for the baseline PF, and \(12\)
Monte Carlo runs per scenario.  The main scenarios are summarized in
Table~\ref{tab:scenarios}.  The final stress case includes missing camera
detections and false detections, and is kept as a transparency check with both
ungated and gated EKF variants.

\begin{table}[htbp]
\centering
\caption{Monte Carlo scenarios used for the filter comparison.}
\label{tab:scenarios}
\begin{tabular}{lcccc}
\toprule
Scenario & Pixel noise [px] & Initial position offset [m] & Dropout & False detections \\
\midrule
Nominal & 8 & \([0.7,-0.5,0.25]\) & 0.00 & 0.00 \\
Poor initialization & 8 & \([2.5,-2.0,0.9]\) & 0.00 & 0.00 \\
High noise & 22 & \([0.9,-0.7,0.25]\) & 0.00 & 0.00 \\
Dropout stress & 12 & \([1.0,-0.8,0.35]\) & 0.18 & 0.06 \\
\bottomrule
\end{tabular}
\end{table}

\subsection{Lie-Group Camera Projection}

For camera \(j\), the world-to-camera extrinsic is
\[
T_{j,cw}
=
\begin{bmatrix}
R_{j,cw} & t_{j,cw}\\
0 & 1
\end{bmatrix}
\in \SE(3).
\]
For a world-frame point \(p_{w,k}\), define the homogeneous point
\[
\bar p_{w,k} =
\begin{bmatrix}
p_{w,k}\\
1
\end{bmatrix}.
\]
The point is transformed into camera coordinates by the \(\SE(3)\) action
\[
\bar p_{j,c,k} = T_{j,cw}\bar p_{w,k},
\qquad
p_{j,c,k} =
\begin{bmatrix}
X_{j,c,k}\\
Y_{j,c,k}\\
Z_{j,c,k}
\end{bmatrix}.
\]
The predicted pixel location is
\[
\hat z_{j,k}
=
h_j(x_k)
=
\pi(K_j,T_{j,cw}\bar p_{w,k}),
\]
where
\[
\pi(K_j,p_{j,c,k})
=
\begin{bmatrix}
f_{x,j}X_{j,c,k}/Z_{j,c,k}+c_{x,j}\\
f_{y,j}Y_{j,c,k}/Z_{j,c,k}+c_{y,j}
\end{bmatrix}.
\]
The measurement model is
\[
z_{j,k}
=
h_j(x_k)+n_{j,k},
\qquad
n_{j,k}\sim \N(0,R_j).
\]

This is the primary Lie-group component of the implemented filters: camera
extrinsics are represented as \(\SE(3)\) elements, and all filters use the same
world-to-camera action before applying the nonlinear projection map familiar
from multiple-view geometry \cite{hartley2004}.

\section{Implemented Filters}

\subsection{Shared Prediction Model}

The implemented Monte Carlo comparison uses a constant-velocity process model:
\[
x_{k+1} =
F_k x_k + q_k,
\qquad
F_k =
\begin{bmatrix}
I_3 & \Delta t I_3\\
0 & I_3
\end{bmatrix}.
\]
Equivalently,
\[
p_{w,k+1}=p_{w,k}+\Delta t\,v_{w,k}+q_{p,k},
\qquad
v_{w,k+1}=v_{w,k}+q_{v,k}.
\]
Because \(x_k\in \R^6\), the prediction state is an additive abelian group.
Writing this portion in Lie-group notation does not change the implemented
equations.  The nontrivial implemented Lie-group component is therefore the
\(\SE(3)\) camera measurement model above.

\subsection{MEKF-Style EKF}

The MEKF-style filter propagates a mean \(\hat x_k\) and covariance \(P_k\).
The predicted covariance is
\[
P^-_k = F_k P_{k-1}F_k^\T + Q_k.
\]
For a valid camera \(j\), the residual is
\[
r_{j,k}=z_{j,k}-h_j(\hat x^-_k).
\]
The projection Jacobian with respect to \(p_w\) is
\[
J_{\pi,j,k} =
\begin{bmatrix}
f_{x,j}/Z & 0 & -f_{x,j}X/Z^2\\
0 & f_{y,j}/Z & -f_{y,j}Y/Z^2
\end{bmatrix},
\]
where \(p_{j,c,k}=[X,Y,Z]^\T\).  Thus
\[
H_{j,k} =
\begin{bmatrix}
J_{\pi,j,k}R_{j,cw} & 0_{2\times 3}
\end{bmatrix}.
\]
Stacking all valid cameras gives \(r_k\), \(H_k\), and \(R_k\).  The correction
is
\[
S_k = H_kP^-_kH_k^\T+R_k,
\qquad
K_k=P^-_kH_k^\T S_k^{-1},
\]
\[
\hat x_k=\hat x^-_k+K_kr_k,
\]
with the Joseph covariance update used for numerical robustness.

\subsection{Iterated EKF and Gated Variants}

The ItEKF uses the same prediction and camera model, but repeatedly
relinearizes \(h(\cdot)\) around the current update estimate.  At iteration
\(\ell\),
\[
\hat z_k^{(\ell)}=h(\hat x_k^{(\ell)}),
\qquad
H_k^{(\ell)}=
\left.\frac{\partial h}{\partial x}\right|_{\hat x_k^{(\ell)}}.
\]
The iterated residual is
\[
\tilde r_k^{(\ell)}
=
z_k-h(\hat x_k^{(\ell)})
+H_k^{(\ell)}(\hat x_k^{(\ell)}-\hat x^-_k),
\]
and the update is
\[
\hat x_k^{(\ell+1)}
=
\hat x^-_k+
K_k^{(\ell)}\tilde r_k^{(\ell)}.
\]
The gated variants reject a camera update when its Mahalanobis distance exceeds
a selected threshold:
\[
r_{j,k}^{\T}S_{j,k}^{-1}r_{j,k}>\gamma.
\]
This is a measurement-validation stage, not a dropout fix: if a camera does not
provide a measurement, both gated and ungated variants simply skip that camera.

\subsection{Bootstrap Particle Filter}

The PF represents the posterior by weighted samples
\[
\mathcal X_k =
\left\{
x_k^{(i)}, w_k^{(i)}
\right\}_{i=1}^{N_p},
\qquad
x_k^{(i)}=
\begin{bmatrix}
p_{w,k}^{(i)}\\
v_{w,k}^{(i)}
\end{bmatrix}.
\]
Particles are predicted through the same constant-velocity model with sampled
process noise.  For each valid camera,
\[
\hat z_{j,k}^{(i)}
=
\pi(K_j,T_{j,cw}\bar p_{w,k}^{(i)}),
\qquad
r_{j,k}^{(i)}
=
z_{j,k}-\hat z_{j,k}^{(i)}.
\]
The particle likelihood is
\[
\ell_k^{(i)}
=
\prod_{j\in\V_k}
\N(r_{j,k}^{(i)};0,R_j),
\]
with a robust likelihood floor in contaminated measurement cases.  The weights
are updated by
\[
w_k^{(i)}
\propto
w_{k|k-1}^{(i)}\ell_k^{(i)},
\qquad
\sum_i w_k^{(i)}=1.
\]
Particles are resampled with probability \(w_k^{(i)}\), and the reported
trajectory estimate is
\[
\hat p_{w,k}
=
\sum_{i=1}^{N_p}w_k^{(i)}p_{w,k}^{(i)}.
\]

\section{Simulation Results}

\subsection{Trajectory and Visual Setup}

Figure~\ref{fig:trajectory} shows the simulated top-view trajectory and camera
layout used for the controlled comparison.  The same geometry is used by all
filters, and all filters receive the same generated detections.

\begin{figure}[htbp]
\centering
\includegraphics[width=0.78\linewidth]{trajectory_top_view.png}
\caption{Top-view simulated trajectory and external-camera layout used in the controlled Monte Carlo comparison.}
\label{fig:trajectory}
\end{figure}

\subsection{Filter Accuracy and Runtime}

Table~\ref{tab:main-results} reports RMSE, 90th percentile error, and mean
runtime per step.  The MEKF and ItEKF have nearly identical accuracy in nominal
and high-noise cases because the posterior remains close to unimodal.  The
ItEKF is more expensive because it repeatedly relinearizes the projection
model.  The PF is slower than the EKFs at the baseline \(450\) particles and is
not the most accurate in easy cases, but it provides a sampled posterior that
becomes valuable when measurements are ambiguous.

\begin{table}[htbp]
\centering
\caption{Main Monte Carlo comparison.  Runtime is mean step time.}
\label{tab:main-results}
\scriptsize
\begin{tabular}{llccc}
\toprule
Scenario & Filter & RMSE [m] & P90 [m] & Runtime [ms] \\
\midrule
Nominal & MEKF & 0.153 & 0.221 & 0.170 \\
Nominal & ItEKF & 0.153 & 0.221 & 0.489 \\
Nominal & MEKF-gated & 0.154 & 0.223 & 0.348 \\
Nominal & ItEKF-gated & 0.154 & 0.222 & 1.017 \\
Nominal & PF & 0.183 & 0.265 & 1.039 \\
\midrule
Poor initialization & MEKF & 0.155 & 0.224 & 0.179 \\
Poor initialization & ItEKF & 0.155 & 0.224 & 0.480 \\
Poor initialization & MEKF-gated & 0.156 & 0.226 & 0.347 \\
Poor initialization & ItEKF-gated & 0.155 & 0.225 & 1.008 \\
Poor initialization & PF & 0.215 & 0.304 & 1.026 \\
\midrule
High noise & MEKF & 0.346 & 0.506 & 0.207 \\
High noise & ItEKF & 0.347 & 0.508 & 0.567 \\
High noise & MEKF-gated & 0.347 & 0.511 & 0.388 \\
High noise & ItEKF-gated & 0.349 & 0.513 & 1.168 \\
High noise & PF & 0.446 & 0.651 & 1.142 \\
\midrule
Dropout stress & MEKF & 0.539 & 0.830 & 0.193 \\
Dropout stress & ItEKF & 0.542 & 0.828 & 0.540 \\
Dropout stress & MEKF-gated & 0.248 & 0.361 & 0.363 \\
Dropout stress & ItEKF-gated & 0.249 & 0.362 & 1.048 \\
Dropout stress & PF & 0.317 & 0.465 & 1.075 \\
\bottomrule
\end{tabular}
\end{table}

\begin{figure}[htbp]
\centering
\begin{minipage}{0.48\linewidth}
\centering
\includegraphics[width=\linewidth]{rmse_bar.png}
\end{minipage}
\hfill
\begin{minipage}{0.48\linewidth}
\centering
\includegraphics[width=\linewidth]{runtime_bar.png}
\end{minipage}
\caption{RMSE and computation-time comparison across implemented filter variants.}
\label{fig:rmse-runtime}
\end{figure}

\begin{figure}[htbp]
\centering
\includegraphics[width=0.88\linewidth]{time_error_comparison.png}
\caption{Representative time history of position error for the simulated comparison.}
\label{fig:time-error}
\end{figure}

\subsection{Particle Count Sensitivity}

The PF improves as the particle count increases, but the improvement saturates
in mostly unimodal cases while runtime grows.  Table~\ref{tab:particle-sweep}
shows the dropout stress sweep.  Increasing from \(450\) to \(14400\) particles
improves RMSE from \(0.317\) m to \(0.287\) m, while runtime rises from
\(0.897\) ms to \(14.055\) ms per step.  This supports the interpretation that
PF should not be claimed as universally superior by simply increasing particle
count; its strength is the ability to represent non-Gaussian uncertainty.

\begin{table}[htbp]
\centering
\caption{PF particle-count sweep in the dropout stress case.}
\label{tab:particle-sweep}
\begin{tabular}{rrrr}
\toprule
Particles & RMSE [m] & P90 [m] & Runtime [ms] \\
\midrule
225 & 0.319 & 0.459 & 0.684 \\
450 & 0.317 & 0.465 & 0.897 \\
900 & 0.296 & 0.434 & 1.319 \\
1800 & 0.292 & 0.431 & 2.065 \\
3600 & 0.289 & 0.419 & 4.167 \\
7200 & 0.288 & 0.419 & 7.747 \\
14400 & 0.287 & 0.415 & 14.055 \\
\bottomrule
\end{tabular}
\end{table}

\begin{figure}[htbp]
\centering
\includegraphics[width=0.72\linewidth]{pf_particle_sweep.png}
\caption{PF particle-count sensitivity.  More particles reduce sampling error but increase runtime.}
\label{fig:particle-sweep}
\end{figure}

\section{A Case Where PF Is Preferable}

The most favorable case for PF is not a simple high-noise Gaussian measurement
model, but an ambiguous visual association model.  Suppose each camera returns
a candidate set
\[
\mathcal Z_{j,k}
=
\{z_{j,k}^{(1)},\ldots,z_{j,k}^{(M_j)}\}.
\]
An EKF-style filter generally needs a single association
\[
\tilde z_{j,k}=z_{j,k}^{(m^\star)}
\]
before applying a Gaussian update.  If \(m^\star\) is wrong, the innovation can
pull the filter toward the wrong target.  The PF can instead use a mixture
likelihood:
\[
L_{j,k}(x_k^{(i)})
=
\epsilon+
\sum_{m=1}^{M_j}
\alpha_m
\N\!\left(
z_{j,k}^{(m)};
h_j(x_k^{(i)}),
R_j
\right),
\]
\[
w_k^{(i)}
\propto
w_{k|k-1}^{(i)}
\prod_j L_{j,k}(x_k^{(i)}).
\]
No candidate label is assigned to a particle.  Instead, each particle is
weighted by how well it explains any plausible candidate.

Table~\ref{tab:ambiguous} shows that this is the regime where PF clearly wins.
The candidate-set PF with \(450\) particles achieves \(0.186\) m RMSE, while
the best gated EKF baseline is \(0.319\) m.  This is the main evidence that PF
is the natural estimator when the posterior is multimodal.

\begin{table}[htbp]
\centering
\caption{Ambiguous candidate-set experiment.}
\label{tab:ambiguous}
\begin{tabular}{lccc}
\toprule
Filter & RMSE [m] & P90 [m] & Runtime [ms] \\
\midrule
MEKF & 0.541 & 0.784 & 0.145 \\
ItEKF & 0.542 & 0.784 & 0.382 \\
MEKF-gated & 0.319 & 0.476 & 0.257 \\
ItEKF-gated & 0.325 & 0.485 & 0.705 \\
PF-candidates & \textbf{0.186} & \textbf{0.266} & 0.969 \\
\bottomrule
\end{tabular}
\end{table}

\begin{figure}[htbp]
\centering
\begin{minipage}{0.48\linewidth}
\centering
\includegraphics[width=\linewidth]{ambiguous_candidate_rmse.png}
\end{minipage}
\hfill
\begin{minipage}{0.48\linewidth}
\centering
\includegraphics[width=\linewidth]{ambiguous_candidate_error.png}
\end{minipage}
\caption{Ambiguous candidate-set case.  The PF uses a mixture likelihood over candidates and maintains multiple hypotheses.}
\label{fig:ambiguous}
\end{figure}

\section{Hardware Calibration Hypothesis}

\subsection{Motivation}

A PURT hardware test on April 24, 2026 produced very large raw PF errors.  The
working hypothesis was that the active-marker rigid body gave a stable QTM
marker pose, but this pose was not the optical camera pose used by the
projection model.  In addition, an old \texttt{camera0.yaml} file described a
\(1600\times1200\) calibration, while the UP camera images used by the tracker
were \(640\times480\).  Either error corrupts the \(\SE(3)\) projection model.

The hypothesis was tested in simulation by replaying an extracted PURT run-1
QTM trajectory and generating clean camera measurements from a correct camera
rig.  The PF was then deliberately run with wrong intrinsics and/or wrong
extrinsics.  This isolates calibration error from detector quality, ROS timing,
and QTM parsing.

\subsection{Calibration Ablation}

The correct UP intrinsics were approximately
\[
K_{\mathrm{UP}}=
\begin{bmatrix}
338.972 & 0 & 320.473\\
0 & 340.546 & 235.618\\
0 & 0 & 1
\end{bmatrix},
\]
whereas the erroneous \texttt{camera0.yaml} intrinsic matrix was
\[
K_{\mathrm{bad}}=
\begin{bmatrix}
847.429 & 0 & 801.181\\
0 & 851.364 & 589.046\\
0 & 0 & 1
\end{bmatrix}.
\]
The wrong-extrinsic cases perturb the assumed camera pose as
\[
\hat R_{j,wc}
=
R_{j,wc}\Exp(\delta\phi_j^\wedge),
\qquad
\hat t_{j,wc}
=
t_{j,wc}+R_{j,wc}\delta c_j,
\]
then convert to world-to-camera form
\[
\hat R_{j,cw}=\hat R_{j,wc}^{\T},
\qquad
\hat t_{j,cw}
=
-\hat R_{j,wc}^{\T}\hat t_{j,wc}.
\]

Table~\ref{tab:hardware} and Fig.~\ref{fig:hardware-ablation} show that the
hypothesis is plausible.  The clean case gives \(0.943\) m PF RMSE.  Using only
\texttt{camera0.yaml} raises the PF RMSE to \(6.654\) m.  A \(15\) deg,
\(0.35\) m marker-like optical-frame error raises it to \(8.692\) m.  Combining
the same marker-like error with \texttt{camera0.yaml} gives \(7.650\) m.  Thus
the calibration failure mode alone is sufficient to reproduce multi-meter
hardware-scale error in an otherwise clean simulator.

\begin{table}[htbp]
\centering
\caption{Hardware calibration hypothesis ablation.}
\label{tab:hardware}
\scriptsize
\begin{tabular}{lccccc}
\toprule
Case & Intrinsics & Extrinsic error & PF RMSE [m] & PF P90 [m] & Pixel P90 [px] \\
\midrule
Correct & UP & correct & 0.943 & 1.351 & 17.1 \\
\texttt{camera0.yaml} only & bad & correct & 6.654 & 9.055 & 632.3 \\
5 deg, 0.10 m & UP & mild & 2.679 & 4.010 & 38.8 \\
15 deg, 0.35 m & UP & marker-like & 8.692 & 15.089 & 112.1 \\
25 deg, 0.70 m & UP & severe & 8.331 & 13.671 & 197.6 \\
15 deg + \texttt{camera0} & bad & marker-like & 7.650 & 10.662 & 738.5 \\
25 deg + \texttt{camera0} & bad & severe & 7.740 & 11.064 & 847.0 \\
\bottomrule
\end{tabular}
\end{table}

\begin{figure}[htbp]
\centering
\begin{minipage}{0.48\linewidth}
\centering
\includegraphics[width=\linewidth]{hardware_calibration_ablation_rmse.png}
\end{minipage}
\hfill
\begin{minipage}{0.48\linewidth}
\centering
\includegraphics[width=\linewidth]{hardware_calibration_pixel_residual.png}
\end{minipage}
\caption{Hardware calibration ablation.  Wrong intrinsics and marker-frame/optical-frame errors produce multi-meter PF errors and large reprojection residuals.}
\label{fig:hardware-ablation}
\end{figure}

\begin{figure}[htbp]
\centering
\includegraphics[width=0.95\linewidth]{hardware_hypothesis_pipeline.png}
\caption{Hardware data in simulation: real ROS-bag snippets, QTM trajectory replay, simulated calibration ablation, and PF error metrics.}
\label{fig:hardware-pipeline}
\end{figure}

\begin{figure}[htbp]
\centering
\begin{minipage}{0.48\linewidth}
\centering
\includegraphics[width=\linewidth]{hardware_purt_camera_layout.png}
\end{minipage}
\hfill
\begin{minipage}{0.38\linewidth}
\centering
\includegraphics[width=\linewidth]{hardware_active_marker_camera.jpeg}
\end{minipage}
\caption{PURT-style camera layout and hardware camera with marker rig.  The ablation models the effect of confusing the marker frame with the optical camera frame.}
\label{fig:hardware-layout}
\end{figure}

\subsection{Practical Hardware Implications}

The hardware result does not prove that calibration was the only issue in the
real run; timing, detector quality, and marker visibility can still contribute.
However, the ablation shows that calibration error is sufficient to explain the
observed error scale.  Before another hardware PF run, QTM marker poses should
be treated as marker-body poses, not optical camera poses; a fixed
\(T_{\mathrm{marker},\mathrm{optical}}\) hand-eye transform should be calibrated
for each camera; and QTM drone positions should be reprojected into each camera
as a pre-filter sanity check.  If the projected point does not land near the
detected drone, the filter should not be run yet.

\section{Proposed Body-Frame-Acceleration Lie-Group PF}

The implemented PF is position/velocity based.  To include more Lie-group
structure in the prediction step, the next extension is to let each particle
carry a latent inertial-navigation state.  This follows the same motivation as
invariant filtering, where body-frame inputs and errors are naturally expressed
on the group \cite{barrau2017}, and connects to particle filtering on Lie
groups \cite{fpf_lg,chahbazian2023}:
\[
X_k^{(i)}
=
\left(
R_{wb,k}^{(i)},v_{w,k}^{(i)},p_{w,k}^{(i)}
\right)
\in \SEtwo.
\]
The output of interest remains the trajectory \(p_{w,k}\), but \(R_{wb,k}\) is
carried so that body-frame IMU acceleration can be mapped into world-frame
motion.

Let the body-frame IMU input be
\[
u_k=(\omega_{m,k},a_{m,k}),
\]
and sample body-frame process noise
\[
\eta_{\omega,k}^{(i)}\sim \N(0,Q_\omega),
\qquad
\eta_{a,k}^{(i)}\sim \N(0,Q_a).
\]
No IMU bias states are included in this proposed version.  The attitude
prediction is
\[
R_{wb,k+1}^{(i)}
=
R_{wb,k}^{(i)}
\Exp_{SO(3)}
\left(
(\omega_{m,k}+\eta_{\omega,k}^{(i)})\Delta t
\right),
\]
where the product is group composition on \(\SO(3)\).  The accelerometer is
mapped into world frame by
\[
a_{w,k}^{(i)}
=
R_{wb,k}^{(i)}
\left(a_{m,k}+\eta_{a,k}^{(i)}\right)
+g_w,
\]
with \(g_w\) omitted if \(a_{m,k}\) is already gravity compensated.  The
velocity and position are then propagated as
\[
v_{w,k+1}^{(i)}
=
v_{w,k}^{(i)}
+\Delta t\,a_{w,k}^{(i)},
\]
\[
p_{w,k+1}^{(i)}
=
p_{w,k}^{(i)}
+\Delta t\,v_{w,k}^{(i)}
+\frac{1}{2}\Delta t^2a_{w,k}^{(i)}.
\]
The camera update remains unchanged:
\[
\hat z_{j,k}^{(i)}
=
\pi(K_j,T_{j,cw}\bar p_{w,k}^{(i)}).
\]
Thus the proposed PF uses Lie groups in both places: \(\SO(3)\) propagation and
body-to-world acceleration mapping in prediction, and \(\SE(3)\) camera
projection in measurement.

\section{Conclusions}

The study shows that Lie-group camera projection is the right mathematical
language for multi-camera UAV tracking: it keeps camera extrinsics explicit and
prevents frame-convention ambiguity.  In the implemented position/velocity
tracking problem, the EKF-style filters are fastest and most accurate in
nominal unimodal cases.  The PF is not the universal winner, but it is the
natural estimator when the posterior is non-Gaussian or multimodal, as shown by
the candidate-set experiment.  The hardware ablation supports the hypothesis
that incorrect intrinsics and marker-frame/optical-frame extrinsic errors can
explain multi-meter hardware PF errors.  The next technical step is a
body-frame-acceleration Lie-group PF in which each particle carries latent
attitude on \(\SO(3)\) or \(\SE_2(3)\), while the reported estimate remains the
UAV trajectory.

\appendix

\section{Projection Jacobian}

Let
\[
p_c=R_{cw}p_w+t_{cw}=
\begin{bmatrix}X&Y&Z\end{bmatrix}^{\T}.
\]
For
\[
u=f_xX/Z+c_x,
\qquad
v=f_yY/Z+c_y,
\]
the projection derivative with respect to the camera-frame point is
\[
J_{\pi}
=
\begin{bmatrix}
f_x/Z & 0 & -f_xX/Z^2\\
0 & f_y/Z & -f_yY/Z^2
\end{bmatrix}.
\]
Since \(\partial p_c/\partial p_w=R_{cw}\), the measurement Jacobian for the
position/velocity state is
\[
H_k
=
\begin{bmatrix}
J_{\pi}R_{cw} & 0_{2\times 3}
\end{bmatrix}.
\]

\section{Iterated EKF Algorithm}

Given \(\hat x^-_k\), \(P^-_k\), and stacked measurement \(z_k\), initialize
\[
\hat x_k^{(0)}=\hat x^-_k.
\]
At each iteration \(\ell\), compute
\[
\hat z_k^{(\ell)}=h(\hat x_k^{(\ell)}),
\qquad
H_k^{(\ell)}
=
\left.\frac{\partial h}{\partial x}\right|_{\hat x_k^{(\ell)}}.
\]
The residual corrected back to the predicted state is
\[
\tilde r_k^{(\ell)}
=
z_k-h(\hat x_k^{(\ell)})
+H_k^{(\ell)}
(\hat x_k^{(\ell)}-\hat x^-_k).
\]
Then
\[
S_k^{(\ell)}
=
H_k^{(\ell)}P^-_kH_k^{(\ell)\T}+R_k,
\]
\[
K_k^{(\ell)}
=
P^-_kH_k^{(\ell)\T}
\left(S_k^{(\ell)}\right)^{-1},
\]
\[
\hat x_k^{(\ell+1)}
=
\hat x^-_k+
K_k^{(\ell)}\tilde r_k^{(\ell)}.
\]
The procedure stops after a fixed number of iterations or when the update
change is below a tolerance.

\section{Left-Invariant Error for the Proposed Lie-Group PF}

The proposed \(\SE_2(3)\) PF does not need a left-invariant error to perform
the particle update.  The PF propagates particles directly and updates weights
by likelihood.  Nevertheless, a left-invariant error is useful for diagnostics
and for describing particle spread on the group.

Given a weighted group mean \(\bar X_k\), define the particle error
\[
\eta_k^{(i)}=\bar X_k^{-1}X_k^{(i)}.
\]
The tangent-space error is
\[
\epsilon_k^{(i)}
=
\Log\left(\bar X_k^{-1}X_k^{(i)}\right)^\vee.
\]
For \(\SE_2(3)\), this contains attitude, velocity, and position error
coordinates:
\[
\epsilon_k^{(i)}
\approx
\begin{bmatrix}
\delta\theta_k^{(i)}\\
\delta v_{b,k}^{(i)}\\
\delta p_{b,k}^{(i)}
\end{bmatrix}.
\]
The \(b\) subscript means that errors are expressed in the body frame of the
reference mean; the stored particle states are still \(v_w^{(i)}\) and
\(p_w^{(i)}\).

The local weighted Lie-group mean minimizes
\[
\bar X_k
=
\arg\min_Y
\frac{1}{2}
\sum_i w_k^{(i)}
\left\|
\Log(Y^{-1}X_k^{(i)})^\vee
\right\|^2.
\]
Let
\[
\epsilon_i=
\Log(\bar X_k^{-1}X_k^{(i)})^\vee.
\]
Perturb the mean by
\[
Y=\bar X_k\Exp(\delta^\wedge).
\]
Then
\[
Y^{-1}X_k^{(i)}
=
\Exp(-\delta^\wedge)\bar X_k^{-1}X_k^{(i)}.
\]
Using the first-order Baker-Campbell-Hausdorff approximation near the mean,
\[
\Log(Y^{-1}X_k^{(i)})^\vee
\approx
\epsilon_i-\delta.
\]
Thus the local cost becomes
\[
J(\delta)
\approx
\frac{1}{2}
\sum_i w_k^{(i)}
\|\epsilon_i-\delta\|^2.
\]
At the optimum, the derivative at \(\delta=0\) must vanish:
\[
\left.\frac{\partial J}{\partial \delta}\right|_{\delta=0}
=
-\sum_iw_k^{(i)}\epsilon_i
=0.
\]
Therefore,
\[
\sum_i w_k^{(i)}
\Log(\bar X_k^{-1}X_k^{(i)})^\vee
\approx 0.
\]
This is the Lie-group analogue of the Euclidean fact that weighted errors
around the mean sum to zero.

\begin{thebibliography}{99}

\bibitem{thrun2005}
S. Thrun, W. Burgard, and D. Fox, \textit{Probabilistic Robotics}, MIT Press,
Cambridge, MA, 2005.

\bibitem{barrau2017}
A. Barrau and S. Bonnabel, ``The Invariant Extended Kalman Filter as a Stable
Observer,'' \textit{IEEE Transactions on Automatic Control}, Vol. 62, No. 4,
2017, pp. 1797--1812.

\bibitem{hartley2004}
R. Hartley and A. Zisserman, \textit{Multiple View Geometry in Computer
Vision}, 2nd ed., Cambridge University Press, Cambridge, UK, 2004.

\bibitem{fpf_lg}
C. Zhang, A. Taghvaei, and P. G. Mehta, ``Feedback Particle Filter on Matrix
Lie Groups,'' arXiv:1701.02416, 2017.

\bibitem{chahbazian2023}
C. Chahbazian, \textit{Particle Filtering on Lie Groups: Application to
Navigation}, Ph.D. thesis, Universite Paris-Saclay, 2023.

\bibitem{yang2022}
Z.-P. Yang, \textit{Unmanned Aerial System Tracking in Urban Canyon
Environments Using External Vision Sensors}, M.S. thesis, Purdue University,
2022.

\end{thebibliography}

\end{document}
